\documentclass[11pt]{article}
\usepackage{amsmath,amssymb,amsthm}
\usepackage[margin=1in]{geometry}
\newtheorem{theorem}{Theorem}
\newtheorem{lemma}[theorem]{Lemma}
\newtheorem{proposition}[theorem]{Proposition}
\newtheorem{conjecture}[theorem]{Conjecture}
\theoremstyle{remark}
\newtheorem{remark}[theorem]{Remark}
\newcommand{\Q}{\mathbb Q}
\newcommand{\Z}{\mathbb Z}
\newcommand{\C}{\mathbb C}
\newcommand{\e}{\mathrm e}
\title{P4: is $U$ differentially algebraic? A negative result about the strategy, not about $U$}
\date{2026-09-26}
\begin{document}
\maketitle

\noindent Labels: PROVED means a complete argument, given the inputs it cites. VERIFIED means
high-precision floating-point numerics, not an interval certificate. HEURISTIC means an argument
without error control. Notation follows \texttt{paper/journal/paper2a.tex} (``paper2a'').

\section*{Summary of outcome}

This note does \emph{not} prove or disprove that $U$ is differentially algebraic (D-algebraic),
i.e.\ that it satisfies some
$P(x,U,U',\ldots,U^{(k)})=0$ for a polynomial $P\not\equiv0$. It gives reasons to distrust the strategy
proposed in \texttt{private/RESEARCH\_PROGRAM\_OPEN.md} row P4 --- ``dense poles on a positive-measure
arc of $|x|=1$ obstructs D-algebraicity, via a Nishioka/Malgrange-type theorem, anchored on
Theorem~\texttt{thm:fredholm}'' --- without exhibiting a decisive counterexample
(\S\ref{sec:counterexample}: an attempted explicit counterexample is shown to be broken on inspection,
post-review; what survives is the qualitative movable-singularities objection of
Remark~\ref{rem:nogeneralize}). It then identifies the correct class of theorems for this question
(\S\ref{sec:realtool}), checks their hypotheses against what paper2a actually proves
(\S\ref{sec:gap}), and finds the single missing ingredient: a \emph{quantitative} lower bound on the
pole-counting function of $U$, which is not established anywhere in this programme. \S\ref{sec:fredholm}
explains why Theorem~\texttt{thm:fredholm} is the wrong anchor: it is a linear (Fredholm/spectral)
statement about a $q$-difference object, not a differential one, and the paper's existing
Nishioka-type results (\texttt{prop:noqdiff}, \texttt{cor:Qminusone}) already exhaust what
difference-Galois methods say here --- they say nothing about $D$-algebraicity. The row stays OPEN.

\section{What is already proved, and what it does not give}\label{sec:proved}

Recall (paper2a \S8, the ``beyond transcendence'' results):
\begin{itemize}
\item $U,V$ are transcendental over $\Q(x)$ (\texttt{thm:U}, \texttt{thm:V}), not $D$-finite
(\texttt{thm:nonDfinite}), satisfy no linear $q$-difference or Mahler equation
(\texttt{prop:noqdiff}), and no $Q=-1$ relation (\texttt{cor:Qminusone}).
\item $|x|=1$ is a natural boundary for $U,V$ (\texttt{thm:natboundary}), by
P\'olya--Bertrandias (\texttt{thm:PB}) applied to the fact that $U,V$ are meromorphic on
$|x|<1$ with infinitely many poles there.
\item The poles of $U,V$ accumulate at \emph{every} point of an explicit arc,
$\arg x\in[\pi/6,5\pi/6]\cup[7\pi/6,11\pi/6]$ (\texttt{cor:arcUV}, from \texttt{thm:arc}). This is a
statement about the closure of the pole set (topological density on the arc); it carries no
information about how many poles lie in $|x|<r$ for a given $r<1$, i.e.\ no counting-function
estimate.
\end{itemize}

The mechanism behind \texttt{thm:nonDfinite} is Lemma~\texttt{lem:odeposes}: a pole in $|x|<1$ of a
solution of a \emph{linear} ODE with coefficients holomorphic on the closed disc must be a zero of
the leading coefficient, hence there are only finitely many such poles unless the leading
coefficient vanishes identically. Infinitely many poles accumulating on the boundary therefore rules
out a linear ODE (with polynomial, or holomorphic-near-$\overline D$, coefficients). This is exactly
why the row's suggested next step reaches for the same kind of pole-count obstruction one dimension
up, for \emph{nonlinear} (algebraic) ODEs.

\begin{remark}[Why this does not just generalize]\label{rem:nogeneralize}
Lemma~\texttt{lem:odeposes} is special to linear equations: the singularities of a solution of a
linear ODE are exactly the \emph{fixed} singularities determined by the equation (zeros of the
leading coefficient), independent of the solution or of initial data. Solutions of \emph{nonlinear}
algebraic ODEs can have \emph{movable} singularities, whose location depends on the solution itself
and need not be confined to a coefficient-determined finite set. Whether the movable singularities
of a given algebraic ODE are poles (rather than branch points or worse) is exactly the content of
Painlev\'e analysis / the Painlev\'e property, and equations that do have the Painlev\'e property
(e.g.\ any order-1 equation $(y')^2=$ polynomial of degree $\le4$ in $y$, the Weierstrass equation)
have solutions whose movable poles can be \emph{dense} on a curve. So ``infinitely many poles dense
on an arc'' is not, by itself, evidence against D-algebraicity; the linear-case argument has no
literal nonlinear analogue.
\end{remark}

\section{No verified counterexample; the real objection is movable singularities}\label{sec:counterexample}

\noindent\textbf{Correction (post-review, Reviewer AH).} An earlier version of this section claimed a
PROVED explicit counterexample: a Weierstrass $\wp$-function pulled back along $\zeta\mapsto
x=\e^{2\pi i\zeta/\omega_2}$ ($\omega_2$ purely imaginary), asserting the sub-lattice $c+n\omega_1$
lands exactly on $|x|=1$ with argument equidistributed by Weyl/Kronecker. This is arithmetically
wrong: with $\omega_1$ real and $\omega_2$ purely imaginary, $\omega_1/\omega_2$ is purely
\emph{imaginary}, not real, so $\zeta/\omega_2=c/\omega_2+n(\omega_1/\omega_2)+m$ has its $n$-dependence
entirely in the \emph{modulus} of $x$ (via $\Im(\zeta/\omega_2)$), not its argument; the sub-lattice
does not stay on $|x|=1$ at all, and shifting by the other generator changes $x$ by a root of unity,
not a dense set. Worse, the two requirements needed for the construction --- $\tau=\omega_2/\omega_1$
non-real (so $\wp$ is a genuine, non-degenerate elliptic function) and $\omega_1/\omega_2$ real (so the
pullback's argument, not modulus, varies along a sub-lattice) --- are mutually exclusive for any
lattice $\Lambda$. Both attempts in the original write-up reduce, on inspection, to the ordinary
$q$-expansion pullback $g(q)=\wp(\zeta)$, $q=\e^{2\pi i\tau}$, whose poles accumulate only at $q=0$ and
$q=\infty$ (the cusps), never on a circle of positive radius --- exactly the ``first attempt'' the
original draft itself flagged as failing, restated with an incorrect modulus computation in the
``fix''. \textbf{No explicit counterexample is exhibited here.} Proposition~\ref{prop:counterexample}
below is downgraded from PROVED to a documented non-result, and Remark~\ref{rem:nogeneralize} --- which
does not depend on this construction --- is the actual load-bearing objection to the ledger's proposed
strategy.

\begin{proposition}[OPEN: no verified example, construction attempted and found broken]\label{prop:counterexample}
We looked for a function $g$, D-algebraic (satisfying a fixed algebraic ODE with polynomial
coefficients), meromorphic on a disc $|x|<1$, with poles dense on a positive-measure arc of $|x|=1$.
We did not find one: the natural candidate (an elliptic function pulled back by an exponential change
of variable along a lattice direction) provably cannot work, by the incompatibility noted above.
Whether such a $g$ exists (via, e.g., an automorphic/Fuchsian-group construction rather than a lattice
pullback, or a genuinely different mechanism) is left open; we do not need it to make the point below,
which rests instead on Remark~\ref{rem:nogeneralize}.
\end{proposition}

\noindent Consequently: \emph{the ledger's strategy is not vindicated by an exhibited counterexample},
but Remark~\ref{rem:nogeneralize} already gives an independent, qualitative reason to distrust it:
Lemma~\texttt{lem:odeposes}'s pole-confinement argument is a fact about \emph{linear} ODEs (fixed
singularities, at zeros of the leading coefficient, independent of the solution); solutions of
\emph{nonlinear} algebraic ODEs can have \emph{movable} singularities whose location depends on
initial data, and when an equation has the Painlev\'e property (e.g.\ the order-1 Weierstrass equation
$(y')^2=4y^3-g_2y-g_3$, or more generally any Painlev\'e-type equation) those movable singularities are
honest poles, not merely branch points, with no a priori constraint tying their location to a
coefficient-determined finite set. This is a structural reason the linear argument does not
literally generalize; it is not by itself a proof that density-on-an-arc is compatible with
D-algebraicity for some specific function, since we have not exhibited one. What it shows is that a
theorem of the form ``dense poles on a positive-measure boundary arc $\Rightarrow$ not D-algebraic''
would have to say something genuinely new about movable singularities of nonlinear ODEs near a natural
boundary -- it is not a corollary of the elementary fixed-singularity fact that makes the linear case
(\texttt{thm:nonDfinite}) work, and no such theorem is currently known to us in citable form. This
matches the ledger's own caveat that ``some Nishioka theorems are about difference equations, not
differential ones'': the density-only line of reasoning is not established to be sufficient, by any
theorem we can cite, regardless of which named result one invokes.

\section{The right class of theorem, and its hypotheses}\label{sec:realtool}

The correct tool for turning growth/pole-density information about a meromorphic function into
non-D-algebraicity is Nevanlinna-theoretic, not the fixed-singularity argument of
\texttt{lem:odeposes}, and not Mahler/Nishioka difference-Galois theory (\S\ref{sec:fredholm}). The
relevant literature is:
\begin{itemize}
\item Malmquist's theorem (1913/1920) and its extensions: a first-order algebraic differential
equation $F(x,y,y')=0$ with a transcendental meromorphic solution on $\C$ of finite order forces the
equation to be of a very restricted (Riccati-like) form.
\item Steinmetz (\emph{Ein Malmquistscher Satz f\"ur algebraische Differentialgleichungen h\"oherer
Ordnung}, J.\ Reine Angew.\ Math.\ 316 (1980)) and Bank--Kaufman (several papers, 1970s-80s, e.g.\
\emph{On the growth of meromorphic solutions of nonlinear differential equations}, Trans.\ AMS 1980)
generalize Malmquist's theorem to higher order and give explicit bounds on the Nevanlinna
characteristic $T(r,y)$ (equivalently, via the first main theorem, on the pole-counting function
$N(r,y)=\int_0^rn(t,y)/t\,dt$) of any transcendental meromorphic solution of a fixed algebraic ODE
with rational-function (or polynomial) coefficients, in terms of the degree of the coefficients and
the order/degree of the equation.
\item The upshot usable here is the contrapositive: if one can show that $n(r,U)$ (the number of
poles of $U$ in $|x|\le r$) grows, as $r\to1^-$, \emph{faster} than any bound compatible with a
solution of \emph{some} algebraic ODE of bounded order and degree, for every order/degree pair, then
$U$ is not D-algebraic.
\end{itemize}

\begin{remark}[the disc vs.\ the plane]\label{rem:disc}
The classical statements (Malmquist, Steinmetz, Bank--Kaufman) are for functions meromorphic on all
of $\C$, where $T(r,y)$ is compared to $T(r)$ of the (polynomial or rational) coefficients as
$r\to\infty$. $U$ is meromorphic only on the disc $|x|<1$, with $|x|=1$ a natural boundary, so the
relevant growth is of $n(r,U)$ or $T(r,U)$ as $r\to1^-$ (necessarily $\to\infty$, since $U$ is not
rational). A full transplant of the Bank--Kaufman/Steinmetz machinery to functions meromorphic only
on a disc, with polynomial ODE coefficients, is not something we could locate a ready-made citation
for; the relevant framework is Tsuji's Nevanlinna theory adapted to the disc (M.\ Tsuji,
\emph{Potential Theory in Modern Function Theory}, 1959, Ch.\ III), and Bank himself has papers
specifically on meromorphic solutions of algebraic differential equations \emph{in the unit disc}
(S.\ Bank, \emph{On the instantaneous singularities of solutions of nonlinear differential
equations}, and related 1970s work on the disc case), but we have not verified in detail that these
give the precise counting-function bound needed here; this is exactly the ``verify applicability
carefully'' caution the programme row asks for, and we flag it honestly as HEURISTIC: the shape
of theorem is right, the citation for the disc case is not one we independently re-derived or
checked line by line.
\end{remark}

\section{Why Theorem \texttt{thm:fredholm} does not anchor this}\label{sec:fredholm}

Theorem~\texttt{thm:fredholm} computes $\det(I-\lambda\mathcal T)$, a Fredholm determinant, in closed
form as a $q$-cosine / Hahn--Exton $q$-Bessel function $F(\lambda,q)=\cos(\sqrt\lambda\,Z;q^2)$. Two
things follow that are relevant here, and neither reaches D-algebraicity of $U(x)$:
\begin{enumerate}
\item It is a statement about $F$ as a function of the spectral parameter $\lambda$ at \emph{fixed}
$q$ (an entire function of $\lambda$, an explicit $q$-Bessel function -- $D$-finite in $\lambda$, in
fact holonomic, since $q$-Bessel functions satisfy linear $q$-difference/differential equations in
their argument). This says nothing about the dependence on $q$ itself, which is where the natural
boundary and the arc-density of \texttt{thm:arc} live.
\item As a function of $q$ (equivalently $x$), $F(1,q)=1-\Sigma_1$ and its zeros are the travel
poles; the operator $\mathcal T=\mathcal T(q)$ satisfies no linear $q$-difference equation in $q$
that we are aware of and none is proved in paper2a -- the ``differential-Galois'' structure available
here is on the \emph{difference} side (Nishioka's theory, as cited at \texttt{prop:noqdiff} for
Mahler equations, and the $Q=-1$ analysis of \texttt{cor:Qminusone}), not a differential-Galois
structure in $x$. Difference-Galois/Nishioka theorems constrain algebraic relations and
$q$-difference or Mahler equations satisfied by $U$; paper2a already proves the strongest available
statements of that kind (no linear $q$-difference or Mahler equation, no $Q=-1$ relation). None of
this is evidence, for or against, an algebraic \emph{differential} equation in $x$: a function can
satisfy no Mahler/difference equation and still be, or fail to be, D-algebraic; the two properties
are logically independent in general (Mahler functions and D-algebraic functions overlap only in
degenerate cases, by the different normal forms Nishioka's and Malmquist/Steinmetz's theorems put on
each).
\end{enumerate}
So thm:fredholm is real machinery for the pole-existence and pole-density results already in the
paper (via the explicit zero set of a $q$-Bessel function), but it is a statement in the wrong
variable (the spectral parameter $\lambda$, or the difference-structure in $q$) to serve as the
anchor for a differential-algebraicity argument in $x$. We could not find, and do not claim to have
found, a route from \texttt{thm:fredholm} to a Nishioka-type differential-algebraic obstruction.

\section{The precise gap, and what would close P4}\label{sec:gap}

Given \S\ref{sec:counterexample}--\ref{sec:realtool}, closing P4 (in either direction) needs:

\begin{enumerate}
\item[(a)] \emph{A quantitative rate.} A lower bound on $n(r,U)$, the number of poles of $U$ in
$|x|\le r$, as $r\to1^-$. \texttt{thm:arc}'s proof gives, near a fixed root of unity $\zeta$ on the
arc, an infinite sequence of poles $q^\ast_j=\zeta\e^{-t^\ast_j}$ with $t^\ast_j\sim t^0_j=(L+i\pi
)/(2j+3/2)$ along an explicit tie ray, i.e.\ $q^\ast_j\to\zeta$ at rate $\sim c/j$. In principle this
already gives, at each of the (dense set of) $\zeta$ on the arc where \texttt{cor:arcUV} applies, a
sequence of poles approaching the boundary at a rate $O(1/j)$ in $t$, which translates to
$1-|q^\ast_j|=O(1/j)$; summing over the countably many $\zeta$ with $N\ge16$ has not been done in
paper2a and does not obviously give a clean closed-form $n(r,U)$ estimate (poles from different
$\zeta$'s interleave in an uncontrolled way; no global counting argument is in the paper). This
summation is the concrete, well-scoped computation that is missing; it looks tractable (each $\zeta$
contributes a $\sim\log(1/(1-r))$-type count near itself, and there are $\sim r^{-1+\epsilon}$-many
admissible $\zeta=e(a/N)$, $N\lesssim1/(1-r)$, on the arc) but was not carried out here, and doing it
rigorously (as opposed to heuristically) would need to control double-counting and the
$N$-dependence of the implied constants in \texttt{thm:arc}'s proof, which are not made explicit
there (see \texttt{rem:arcUV}: ``the index beyond which $q^\ast_j$ is a pole $\ldots$ is not made
explicit'').
\item[(b)] \emph{The disc-analog counting theorem.} A Bank--Kaufman/Steinmetz-type upper bound, valid
for meromorphic functions on a disc with a natural boundary, showing that no algebraic ODE of order
$k$ and degree $d$ (for any fixed $k,d$) can have a solution meromorphic on $|x|<1$ whose
pole-counting function $n(r,\cdot)$ grows faster than [some explicit function of $r$, $k$, $d$] as
$r\to1^-$. We have identified the right family of theorems (\S\ref{sec:realtool}) but have not
verified that a version with the precise hypotheses needed (arbitrary $k,d$, natural boundary rather
than entire coefficients) exists in the literature in citable form; \texttt{rem:disc} states this
honestly as unresolved.
\item[(c)] Compare (a) and (b) for a contradiction, for every $k,d$ simultaneously (or via a
uniform-in-$(k,d)$ argument, e.g.\ showing $n(r,U)$ grows faster than any bound of the Bank--Kaufman
shape, which typically depends polynomially or via a fixed exponential rate on $k,d$ once $r\to1$,
whereas one would need superexponential growth in $1/(1-r)$ uniformly to beat all $(k,d)$ at once,
or a separate argument bounding the ODE's order/degree first, e.g.\ by degree considerations on
residues as in \texttt{cor:Qminusone}'s use of $2$-adic valuations).
\end{enumerate}

None of (a), (b), (c) is done. \textbf{P4 is not closed by this note.} What this note contributes is
negative-but-useful: it rules out the literal strategy stated in the ledger row (density on an arc
$\Rightarrow$ non-D-algebraic, via thm:fredholm), replacing it with the correct
(Nevanlinna/Bank--Kaufman/Steinmetz) family of theorems and pinpointing the exact missing input, a
quantitative pole-counting rate together with a disc-version growth theorem, so that the next
attempt does not repeat the invalid step of \S\ref{sec:counterexample}.

\section*{Weakest points}
\begin{itemize}
\item \S\ref{sec:counterexample}'s original elliptic-pullback construction is \textbf{wrong}, not
merely under-polished: Reviewer AH caught that the claim ``$|x|=1$ exactly, since $\omega_1$ is real''
(original line 132) is arithmetically false for the stated map, and that the underlying requirements
($\tau=\omega_2/\omega_1$ non-real for $\wp$ to be non-degenerate, vs.\ $\omega_1/\omega_2$ real for the
pullback's argument to vary along $|x|=1$) are mutually exclusive. This is fixed above by downgrading
Proposition~\ref{prop:counterexample} to an open, unresolved question and relying instead on the
qualitative Remark~\ref{rem:nogeneralize} (movable singularities / Painlev\'e property), which does not
depend on the broken construction. A genuine explicit example (if one exists) would need a different
mechanism, e.g.\ Fuchsian/automorphic rather than lattice-elliptic.
\item \S\ref{sec:realtool}/\ref{rem:disc}: we cite Malmquist, Steinmetz, and Bank--Kaufman from
memory of the theory's shape (Nevanlinna-theoretic growth bounds for meromorphic solutions of
algebraic ODEs); we did not independently re-derive or numerically check any of these statements in
this session, and we explicitly flag that the disc/natural-boundary version of the theory is the
part we could not pin down a precise citable statement for. Treat the bibliographic pointers here as
HEURISTIC memory, to be checked against the primary sources before citing in a paper.
\item \S\ref{sec:gap}(a)'s claimed shape of a global counting estimate ($n(r,U)\sim$ combining
$\sim r^{-1}$ many arcs each contributing a logarithmic count) is a back-of-envelope estimate, not a
proof; it is plausible from \texttt{thm:arc}'s local rate but we did not attempt the summation.
\item No computation (Rust or otherwise) was run for this note: the question, as scoped, turned out
to be about which theorem applies and why, not about a numerical question the existing tools (t1series,
u5b, etc.) could resolve. If a future pass pursues (a), a Rust enumeration of travel-pole counts
$n(r,U)$ at a sequence of $r\to1^-$, cross-checked against the explicit $t^\ast_j$ formula, would be
the natural first computational step, wrapped under the standing \texttt{perl -e 'alarm 290; exec
@ARGV'} / $\le3$GB / fresh-\texttt{sysctl} discipline.
\end{itemize}

\end{document}
